\chapter{Bảng tổng hợp kết quả bổ sung}
\label{app:per_video_results}

Phụ lục này tổng hợp lại các kết quả F1 được dùng trong phần thực nghiệm và thảo luận của khóa luận. Hai bảng đầu nhắc lại các kết quả rút gọn trên SHOT và ClipShots; bảng cuối cung cấp danh sách đầy đủ các nhóm thí nghiệm, bao gồm cả các dòng dùng cho phần phân tích ở các chương trước.

\section{Bộ dữ liệu SHOT}
\label{app:shot_detail}

Bảng~\ref{tab:per_video_shot} nhắc lại kết quả trung bình của ba mô hình tiêu biểu trên tập kiểm tra SHOT.

\begin{table}[H]
\centering
\small
\renewcommand{\arraystretch}{1.15}
\caption{Kết quả F1 trung bình trên tập kiểm tra của SHOT}
\label{tab:per_video_shot}
\begin{tabular}{L{3.8cm}C{2.9cm}C{2.9cm}C{2.9cm}}
\toprule
\textbf{Tập dữ liệu} & \textbf{TransNetV2} & \textbf{AutoShot} & \textbf{AutoShotV2} \\
\midrule
\AppendixShotRows
\bottomrule
\end{tabular}
\end{table}

\noindent\textbf{Ghi chú đọc bảng:}
\begin{itemize}
    \item Các giá trị trong bảng khớp với Bảng~\ref{tab:autoshotv2_results} trong Chương~\ref{ch:ket_qua_thuc_nghiem}.
    \item AutoShotV2 đạt F1 cao nhất trên SHOT, phù hợp với mục tiêu tối ưu cho miền video ngắn.
\end{itemize}

\section{Bộ dữ liệu ClipShots}
\label{app:clipshots_detail}

Bảng~\ref{tab:per_video_clipshots} tổng hợp kết quả F1 trên ClipShots. Đáng chú ý, ClipShots có miền nội dung đa dạng hơn SHOT, do đó được dùng để quan sát khả năng tổng quát hóa của các biến thể.

\begin{table}[H]
\centering
\small
\renewcommand{\arraystretch}{1.15}
\caption{Kết quả F1 trung bình trên ClipShots}
\label{tab:per_video_clipshots}
\begin{tabular}{L{3.8cm}C{2.9cm}C{2.9cm}C{2.9cm}}
\toprule
\textbf{Tập dữ liệu} & \textbf{TransNetV2} & \textbf{AutoShot} & \textbf{AutoShotV2} \\
\midrule
\AppendixClipshotsRows
\bottomrule
\end{tabular}
\end{table}

\noindent\textbf{Quan sát đáng chú ý.}
Mức suy giảm khoảng $3.4\%$ của AutoShotV2 so với AutoShot theo báo cáo gốc trên ClipShots cho thấy cấu hình tối ưu cho SHOT chưa chuyển dịch ổn định sang miền dữ liệu đa dạng hơn. Phân tích sâu hơn (xem Chương~\ref{ch:ket_luan}, Mục~\ref{sec:limitations}) chỉ ra một số hướng giải thích cần kiểm chứng:
\begin{enumerate}
    \item Hệ số temperature scaling $T^*$ và ngưỡng quyết định được tối ưu trên SHOT có thể bị lệch phân phối khi áp dụng sang ClipShots.
    \item Nhãn many-hot cố định chưa mô tả tốt các chuyển cảnh có độ dài khác nhau, đặc biệt là nhóm chuyển cảnh dần.
    \item Cấu hình Focal Loss chưa thể hiện lợi ích rõ ràng trong ma trận đối chứng hiện tại, nên cần được kiểm chứng thêm trên từng miền dữ liệu.
\end{enumerate}

Các bảng ở hai mục đầu của phụ lục này chỉ đóng vai trò tổng hợp lại số liệu đã dùng trong phần thực nghiệm chính trên SHOT và ClipShots. Việc kiểm chứng định lượng sâu hơn theo từng video cần được thực hiện khi có đầy đủ \textit{log} hoặc CSV chi tiết cho từng mẫu kiểm tra.

\section{Bảng đầy đủ tất cả thực nghiệm}
\label{app:all_experiments}

Bảng~\ref{tab:all_experiments_master} tổng hợp tất cả các dòng thực nghiệm đã được sử dụng trong Chương~\ref{ch:ket_qua_thuc_nghiem} và phần thảo luận ở Chương~\ref{ch:ket_luan}. Các giá trị trong ba cột bộ dữ liệu là F1-score. Do các nhóm thí nghiệm dùng protocol khác nhau, bảng này chỉ dùng để tra cứu tổng quan; các kết luận chính cần đọc theo từng nhóm thí nghiệm ở các chương trước.

{\small
\setlength{\tabcolsep}{2pt}
\renewcommand{\arraystretch}{1.12}
\begin{longtable}{L{2.7cm}L{4.2cm}C{1.35cm}C{1.35cm}C{1.65cm}L{3.3cm}}
\caption[Bảng đầy đủ tất cả thực nghiệm]{Bảng đầy đủ tất cả thực nghiệm. Các giá trị SHOT, BBC và ClipShots là F1-score.}\label{tab:all_experiments_master}\\
\toprule
\textbf{Nhóm} & \textbf{Thực nghiệm / protocol} & \textbf{SHOT} & \textbf{BBC} & \textbf{ClipShots} & \textbf{Ghi chú} \\
\midrule
\endfirsthead
\caption[]{Bảng đầy đủ tất cả thực nghiệm (tiếp theo)}\\
\toprule
\textbf{Nhóm} & \textbf{Thực nghiệm / protocol} & \textbf{SHOT} & \textbf{BBC} & \textbf{ClipShots} & \textbf{Ghi chú} \\
\midrule
\endhead
\midrule
\multicolumn{6}{r}{\textit{Còn tiếp ở trang sau}} \\
\endfoot
\bottomrule
\endlastfoot

\AllExperimentsRows
\end{longtable}
}





